\clearpage
\section{Product Backlog và Definition of Done}

\subsection{Bảng Product Backlog}
\label{sec:product-backlog}

Các task được sắp xếp từ dễ đến khó và được lên kế hoạch giải quyết toàn bộ trong 1 Sprint kéo dài 10 ngày.
Story Points (SP) là thang đo công sức tương đối; cách ước lượng được trình bày
tại Mục~\ref{sec:story-points}.

\newcommand{\code}[1]{\texttt{#1}}
\setlist[itemize]{nosep,leftmargin=*,topsep=0pt,partopsep=0pt}

\renewcommand{\arraystretch}{1.25}
\small
\begin{longtable}{@{}>{\centering\arraybackslash}p{0.5cm}
                   >{\raggedright\arraybackslash}p{2.8cm}
                   >{\raggedright\arraybackslash}p{7.8cm}
                   >{\centering\arraybackslash}p{0.7cm}
                   >{\centering\arraybackslash}p{1.0cm}@{}}
\caption{Product Backlog BrewLite gồm 10 task}\label{tab:product-backlog}\\
\toprule
\textbf{\#} & \textbf{Task (Story)} & \textbf{Yêu cầu / Tiêu chí chấp nhận chính} &
\textbf{SP} & \textbf{Sprint} \\
\midrule
\endfirsthead
\toprule
\textbf{\#} & \textbf{Task (Story)} & \textbf{Yêu cầu / Tiêu chí chấp nhận chính} &
\textbf{SP} & \textbf{Sprint} \\
\midrule
\endhead
\endfoot
\bottomrule
\endlastfoot

1 & Khởi tạo dự án và cấu trúc &
\begin{itemize}
  \item Monorepo hoặc 2 thư mục \code{frontend} (Next.js) và \code{backend} (NestJS).
  \item Cả hai chạy được và trả về ``Hello''.
  \item Có README, Git, \code{.env.example}.
\end{itemize} & 1 & 1 \\
\midrule

2 & API danh sách sản phẩm &
\begin{itemize}
  \item \code{GET /products} trả JSON menu gồm \code{id}, tên, giá, ảnh.
  \item Dữ liệu lấy từ dữ liệu mẫu hoặc CSDL.
  \item \code{GET /products/:id} trả chi tiết sản phẩm tương ứng; mã không tồn tại được xử lý là không tìm thấy.
\end{itemize} & 1 & 1 \\
\midrule

3 & Trang Menu (Frontend) &
\begin{itemize}
  \item Next.js gọi \code{/products} và render lưới sản phẩm (hình, tên, giá).
  \item Có trạng thái \emph{loading} và \emph{empty}.
\end{itemize} & 2 & 1 \\
\midrule

4 & Chi tiết và tùy chọn sản phẩm &
\begin{itemize}
  \item Màn hình chi tiết cho phép chọn size (S/M/L) và topping.
  \item Giá được tính lại theo tùy chọn đã chọn.
\end{itemize} & 2 & 1 \\
\midrule

5 & Giỏ hàng (Cart) &
\begin{itemize}
  \item Thêm, sửa số lượng, xóa sản phẩm trong giỏ.
  \item Lưu state bằng Zustand hoặc Context.
  \item Tự tính tổng tiền; hiển thị badge số lượng.
\end{itemize} & 2 & 1 \\
\midrule

6 & API tạo đơn hàng &
\begin{itemize}
  \item \code{POST /orders} nhận giỏ hàng và validate bằng \code{class-validator}.
  \item Lưu đơn ở trạng thái \code{PENDING}, trả về mã đơn.
\end{itemize} & 2 & 1 \\
\midrule

7 & Đăng ký / Đăng nhập (JWT) &
\begin{itemize}
  \item \code{POST /auth/register} và \code{POST /auth/login}.
  \item Mật khẩu băm bằng bcrypt; phát JWT.
  \item Guard bảo vệ route đặt đơn.
\end{itemize} & 2 & 1 \\
\midrule

8 & Thanh toán không tiền mặt &
\begin{itemize}
  \item Tích hợp mock payment: chọn Ví hoặc Thẻ, gọi \code{POST /payments}.
  \item Thành công: đơn chuyển \code{PAID}.
  \item Lỗi: đơn chuyển \code{PAYMENT\_FAILED}, giỏ hàng được giữ nguyên.
\end{itemize} & 2 & 1 \\
\midrule

9 & Xác nhận, lịch sử đơn và bàn giao &
\begin{itemize}
  \item Màn hình xác nhận hiển thị mã đơn và trạng thái.
  \item \code{GET /orders/me} trả lịch sử đơn của user.
  \item \code{docker-compose} chạy cả 3 dịch vụ (frontend, backend, DB).
  \item Có README và demo end-to-end.
\end{itemize} & 3 & 1 \\
\midrule

10 & Nghiệp vụ backend (Xử lý đặt hàng và thanh toán) &
\begin{itemize}
  \item \textbf{State Machine} đơn hàng: \code{PENDING} $\to$ \code{PAID} $\to$
        \code{PREPARING} $\to$ \code{READY} $\to$ \code{COMPLETED}; chặn chuyển
        trạng thái không hợp lệ.
  \item \textbf{Thanh toán idempotent} bằng \code{Idempotency-Key}.
  \item \textbf{Kiểm soát tồn kho} khi thanh toán bằng transaction và
        optimistic locking theo \code{version} và \code{stock}.
  \item \textbf{Khuyến mãi và điểm thưởng}: áp mã hợp lệ, cộng điểm sau khi đơn \code{PAID}.
  \item Có unit/integration test chứng minh: (a) chặn chuyển trạng thái sai;
        (b) yêu cầu lặp với cùng \code{Idempotency-Key} không xử lý thanh toán lần hai;
        (c) mô phỏng hai khách mua đồng thời để kiểm tra không vượt tồn kho.
\end{itemize} & 3 & 1 \\
\midrule
\multicolumn{3}{r}{\textbf{Tổng SP}} & \textbf{20} & \\
\end{longtable}
\normalsize

\subsection{Phân bổ Story Points}
\label{sec:story-points}

\subsubsection{Nguyên tắc ước lượng}

Nhóm dùng thang Fibonacci rút gọn (\textbf{1, 2, 3, 5, 8, 13}) và ước lượng theo
phương pháp \emph{Planning Poker} trong buổi Sprint Planning. Story Points đo
\emph{công sức tương đối}, không đo giờ làm việc. Mỗi task được so sánh với một
task chuẩn: \textbf{Task~2 (API danh sách sản phẩm) = 1 SP}.

Mỗi task được cân nhắc theo bốn yếu tố:
\begin{itemize}
  \item \textbf{Độ phức tạp kỹ thuật}: mức độ khó của logic, thuật toán, tích hợp.
  \item \textbf{Khối lượng công việc}: số màn hình, số endpoint, số thực thể cần làm.
  \item \textbf{Rủi ro và độ bất định}: công nghệ mới với nhóm, yêu cầu chưa rõ.
  \item \textbf{Phụ thuộc}: task có phải chờ hoặc ảnh hưởng đến task khác hay không.
\end{itemize}

\subsubsection{Giải thích từng hạng mục}

\small
\begin{longtable}{@{}>{\centering\arraybackslash}p{0.6cm}
                   >{\raggedright\arraybackslash}p{3.0cm}
                   >{\centering\arraybackslash}p{0.8cm}
                   >{\raggedright\arraybackslash}p{9.2cm}@{}}
\caption{Giải thích Story Points cho từng task}\label{tab:story-points-estimation}\\
\toprule
\textbf{\#} & \textbf{Task} & \textbf{SP} & \textbf{Lý do ước lượng} \\
\midrule
\endfirsthead
\toprule
\textbf{\#} & \textbf{Task} & \textbf{SP} & \textbf{Lý do ước lượng} \\
\midrule
\endhead
\bottomrule
\endlastfoot

1 & Khởi tạo dự án & 1 &
Khởi tạo cấu trúc và cấu hình nền tảng tương đối đơn giản, ít nghiệp vụ và ít rủi ro. \\
\addlinespace
2 & API danh sách sản phẩm & 1 &
\textbf{Task chuẩn.} Một endpoint đọc danh sách sản phẩm từ CSDL, không có
nghiệp vụ phức tạp. \\
\addlinespace
3 & Trang Menu & 2 &
Gọi API, render lưới và xử lý loading/empty. Nhiều hơn Task~2 vì có UI, state bất đồng bộ
và phụ thuộc vào Task~2. \\
\addlinespace
4 & Chi tiết và tùy chọn & 2 &
Có logic tính giá theo size/topping và phối hợp nhiều thành phần giao diện. Quy tắc
giá cần được PO xác nhận. \\
\addlinespace
5 & Giỏ hàng & 2 &
Quản lý state toàn cục, nhiều thao tác (thêm, sửa, xóa), tính tổng tiền và badge.
Logic thuần frontend, rủi ro thấp. \\
\addlinespace
6 & API tạo đơn hàng & 2 &
Phải thiết kế Order/OrderItem, kiểm tra dữ liệu đầu vào, lưu nhiều bảng trong một
transaction và tính lại tổng tiền ở server. Ảnh hưởng trực tiếp đến Task 8, 9, 10. \\
\addlinespace
7 & Đăng ký / Đăng nhập (JWT) & 2 &
Liên quan bảo mật: băm mật khẩu, phát và xác thực JWT, guard (Passport). Cần phối hợp
backend với frontend và có rủi ro lỗi bảo mật nếu làm sai. \\
\addlinespace
8 & Thanh toán không tiền mặt & 2 &
Tích hợp dịch vụ thanh toán giả lập, xử lý hai nhánh thành công/lỗi, cập nhật trạng thái
đơn và giữ giỏ hàng khi lỗi. Phụ thuộc Task 6 và 7. \\
\addlinespace
9 & Xác nhận, lịch sử đơn và bàn giao & 3 &
Gồm màn hình, endpoint lịch sử và đóng gói \code{docker-compose} cho 3 dịch vụ. Phần
bàn giao dễ phát sinh lỗi môi trường. \\
\addlinespace
10 & Nghiệp vụ backend & 3 &
Gồm state machine, idempotency, kiểm soát tồn kho và khuyến mãi/điểm thưởng. Đây là
nhóm nghiệp vụ phức tạp hơn các task giao diện cơ bản; phạm vi được giới hạn bằng
các ca kiểm thử đã mô tả trong repo. \\
\midrule
\multicolumn{2}{r}{\textbf{Tổng}} & \textbf{20} & \\
\end{longtable}
\normalsize

\subsection{Definition of Done (DoD)}
\label{sec:dod}

Một task được xem là \textbf{hoàn thành} khi thỏa
\emph{tất cả} điều kiện dưới đây. Điều kiện nào chưa đạt thì task vẫn ở trạng thái
``đang làm''.

\subsubsection{DoD chung cho mọi task}

\begin{table}[H]
\centering
\caption{Definition of Done chung cho các task}
\label{tab:dod-checklist}
\renewcommand{\arraystretch}{1.3}
\small
\begin{tabular}{@{}>{\centering\arraybackslash}p{0.6cm}
                 >{\raggedright\arraybackslash}p{3.2cm}
                 >{\raggedright\arraybackslash}p{9.6cm}@{}}
\toprule
\textbf{\#} & \textbf{Nhóm} & \textbf{Điều kiện hoàn thành} \\
\midrule
1 & Code và build &
Code chạy được, không lỗi build; đáp ứng đầy đủ tiêu chí chấp nhận của task. \\
2 & Git và review &
Có commit trên Git với mô tả rõ ràng; code được review trong nhóm trước khi merge. \\
3 & API &
Endpoint có validate đầu vào (\code{class-validator}); trả mã lỗi và thông báo hợp lý. \\
4 & Giao diện &
UI xử lý trạng thái loading và lỗi cơ bản. \\
5 & Tài liệu &
Có hướng dẫn chạy trong README. \\
\bottomrule
\end{tabular}
\end{table}

\subsubsection{DoD bổ sung (áp dụng theo loại task)}

\begin{itemize}
  \item \textbf{Task backend (2, 6, 7, 8, 9, 10):} API xử lý đầu vào hợp lệ và không hợp lệ;
        mật khẩu được băm, route cần bảo vệ có JWT guard.
  \item \textbf{Task frontend (3, 4, 5, 9):} giao diện hiển thị trạng thái tải/lỗi phù hợp
        với luồng của task và cho phép hoàn tất kịch bản người dùng tương ứng.
  \item \textbf{Task 9 (bàn giao):} \code{docker-compose up} chạy được cả frontend,
        backend và DB; demo end-to-end thành công.
  \item \textbf{Task 10 (nghiệp vụ backend):} có kiểm thử tự động cho
        (a) chặn chuyển trạng thái sai, (b) 2 lần cùng \code{Idempotency-Key} chỉ
        tạo 1 đơn, (c) đặt đồng thời không vượt tồn kho; tất cả test pass.
\end{itemize}

\subsubsection{Danh mục kiểm tra trước khi nghiệm thu}

Trước mỗi Sprint Review, nhóm dùng danh mục sau để tự kiểm tra từng task:

\begin{itemize}
  \item[\fbox{\rule{0pt}{1.1ex}\rule{1.1ex}{0pt}}] Code chạy được, không lỗi build.
  \item[\fbox{\rule{0pt}{1.1ex}\rule{1.1ex}{0pt}}] Đáp ứng đủ tiêu chí chấp nhận trong Product Backlog.
  \item[\fbox{\rule{0pt}{1.1ex}\rule{1.1ex}{0pt}}] Có commit Git với mô tả rõ ràng.
  \item[\fbox{\rule{0pt}{1.1ex}\rule{1.1ex}{0pt}}] Đã được review bởi ít nhất một thành viên khác.
  \item[\fbox{\rule{0pt}{1.1ex}\rule{1.1ex}{0pt}}] API có validate đầu vào.
  \item[\fbox{\rule{0pt}{1.1ex}\rule{1.1ex}{0pt}}] UI xử lý loading/lỗi cơ bản.
  \item[\fbox{\rule{0pt}{1.1ex}\rule{1.1ex}{0pt}}] README có hướng dẫn chạy và được cập nhật.
  \item[\fbox{\rule{0pt}{1.1ex}\rule{1.1ex}{0pt}}] (Task 10) Test (a), (b), (c) đều pass.
\end{itemize}
